%% Generated by Sphinx.
\def\sphinxdocclass{report}
\documentclass[letterpaper,10pt,brazil]{sphinxmanual}
\ifdefined\pdfpxdimen
   \let\sphinxpxdimen\pdfpxdimen\else\newdimen\sphinxpxdimen
\fi \sphinxpxdimen=.75bp\relax
\ifdefined\pdfimageresolution
    \pdfimageresolution= \numexpr \dimexpr1in\relax/\sphinxpxdimen\relax
\fi
%% let collapsible pdf bookmarks panel have high depth per default
\PassOptionsToPackage{bookmarksdepth=5}{hyperref}

\PassOptionsToPackage{booktabs}{sphinx}
\PassOptionsToPackage{colorrows}{sphinx}

\PassOptionsToPackage{warn}{textcomp}
\usepackage[utf8]{inputenc}
\ifdefined\DeclareUnicodeCharacter
% support both utf8 and utf8x syntaxes
  \ifdefined\DeclareUnicodeCharacterAsOptional
    \def\sphinxDUC#1{\DeclareUnicodeCharacter{"#1}}
  \else
    \let\sphinxDUC\DeclareUnicodeCharacter
  \fi
  \sphinxDUC{00A0}{\nobreakspace}
  \sphinxDUC{2500}{\sphinxunichar{2500}}
  \sphinxDUC{2502}{\sphinxunichar{2502}}
  \sphinxDUC{2514}{\sphinxunichar{2514}}
  \sphinxDUC{251C}{\sphinxunichar{251C}}
  \sphinxDUC{2572}{\textbackslash}
\fi
\usepackage{cmap}
\usepackage[T1]{fontenc}
\usepackage{amsmath,amssymb,amstext}
\usepackage{babel}



\usepackage{tgtermes}
\usepackage{tgheros}
\renewcommand{\ttdefault}{txtt}



\usepackage[Sonny]{fncychap}
\ChNameVar{\Large\normalfont\sffamily}
\ChTitleVar{\Large\normalfont\sffamily}
\usepackage{sphinx}

\fvset{fontsize=auto}
\usepackage{geometry}


% Include hyperref last.
\usepackage{hyperref}
% Fix anchor placement for figures with captions.
\usepackage{hypcap}% it must be loaded after hyperref.
% Set up styles of URL: it should be placed after hyperref.
\urlstyle{same}

\addto\captionsbrazil{\renewcommand{\contentsname}{Contents:}}

\usepackage{sphinxmessages}
\setcounter{tocdepth}{1}



\title{Lógica de programação com Jogos}
\date{16 ago. 2025}
\release{1.0}
\author{Hamilton V.\@{} Rodrigues}
\newcommand{\sphinxlogo}{\vbox{}}
\renewcommand{\releasename}{Release}
\makeindex
\begin{document}

\ifdefined\shorthandoff
  \ifnum\catcode`\=\string=\active\shorthandoff{=}\fi
  \ifnum\catcode`\"=\active\shorthandoff{"}\fi
\fi

\pagestyle{empty}
\sphinxmaketitle
\pagestyle{plain}
\sphinxtableofcontents
\pagestyle{normal}
\phantomsection\label{\detokenize{index::doc}}

\subsubsection*{Introdução}

\sphinxAtStartPar
O Curso \sphinxstyleemphasis{introdução à lógica de programação}, reúne um compêndio de exercícios e
explicações com ferramentas conhecidas que podem facilitar a criança, adolescente ou adulto na
introdução ao mundo da programação. Possuímos atualmente a versão web e PDF, caso tenha interesse.

\sphinxAtStartPar
Como está dividido esse documento?
\begin{itemize}
\item {} 
\sphinxAtStartPar
Navegação ao Sistema Operacional

\item {} 
\sphinxAtStartPar
Comandos CLI no Sistema Operacional

\item {} 
\sphinxAtStartPar
REGEX

\end{itemize}
\subsubsection*{Engine Defold}

\sphinxAtStartPar
Engine para desenvolvimento de jogos, focada em desenvolvimento de jogos de forma visual.
\subsubsection*{Carga horária}

\sphinxAtStartPar
As aulas de Programação com jogos totalizam 20 horas, sendo uma hora e meia a cada
encontro, duas vezes durante a semana, totalizando 13 aulas em 2 meses.
\subsubsection*{Objetivos Gerais}

\sphinxAtStartPar
O curso de programação com games busca desenvolver nos alunos o pensamento
computacional e o raciocínio lógico, de maneira a resolver problemas mais complexos de
forma lúdica.
\subsubsection*{Objetivos Específicos}

\sphinxAtStartPar
Operações básicas no Sistema Operacional:
Aprender a utilizar o básico do sistema operacional, ajudando a introduzir os alunos à
informática básica, promovendo um melhor aprendizado nas ferramentas propostas
no curso (que exigem conhecimento básico).

\sphinxAtStartPar
Aprender a desenvolver soluções na linguagem \sphinxhref{https://www.lua.org/}{lua} .

\sphinxAtStartPar
Os princípios essenciais a serem aprendidos: condições, laços, variáveis,
acumuladores, eventos, transmissão de informação, edição de imagens dos
personagens. Utilização livre e espontânea da ferramenta.

\sphinxAtStartPar
Como princípios essenciais a serem desenvolvidos pelos alunos está, em primeiro
lugar, entender o funcionamento básico da ferramenta: navegação em tela, salvar e
abrir arquivos, criação de pastas no Windows, salvar imagens do Google no
computador pessoal. As ações poderão ser desenvolvidas utilizando como base os
princípios do pensamento lógico. No decorrer das aulas o aluno aprenderá a criar
personagens, cenários, mudança de cenário, disparos de projéteis, criação de
barras de vida e contadores de moedas. Durante as aulas os alunos aprenderão a
ativar animações em determinadas situações, trabalhando com eventos.

\sphinxstepscope


\chapter{Introdução ao Curso}
\label{\detokenize{modelo_aula_01:introducao-ao-curso}}\label{\detokenize{modelo_aula_01::doc}}

\section{Introdução a lógica de programação com jogos}
\label{\detokenize{modelo_aula_01:introducao-a-logica-de-programacao-com-jogos}}
\sphinxAtStartPar
\sphinxstylestrong{Tempo de Aula:} 1h30min
\subsubsection*{Recursos Didáticos}
\begin{itemize}
\item {} 
\sphinxAtStartPar
Projetor: Exibição de exemplos, Slides e exercícios

\item {} 
\sphinxAtStartPar
Quadro branco: Exibição de mecânicas de forma desenhada

\item {} 
\sphinxAtStartPar
Computadores: Computadores com os softwares citados a baixo instaldos e internet.

\end{itemize}
\subsubsection*{Metodologia}
\begin{itemize}
\item {} 
\sphinxAtStartPar
Expositiva, aplicação de exercícios e solução de problemas.

\item {} 
\sphinxAtStartPar
Exercícios Guiados em programação.

\item {} 
\sphinxAtStartPar
Atividade de exploração Criativa.

\item {} 
\sphinxAtStartPar
Estimulo a participação.

\end{itemize}
\subsubsection*{Objetivos}
\begin{itemize}
\item {} 
\sphinxAtStartPar
Compreender conceitos básicos que estão por trás da criação de jogos.

\item {} 
\sphinxAtStartPar
Conhecer e ter uma visão Geral das ferramentas gráficas para criação de sprites e artes digitais.

\item {} 
\sphinxAtStartPar
Desenvolver raciocínio lógico e habilidades criativas por meio de exercícios práticos.

\end{itemize}
\subsubsection*{Conteúdo Programático}

\sphinxAtStartPar
\sphinxstylestrong{Ambientação}
\begin{itemize}
\item {} 
\sphinxAtStartPar
Ligar computador.

\item {} 
\sphinxAtStartPar
Pesquisar conteúdos no navegador web.

\item {} 
\sphinxAtStartPar
Criação de Pastas.

\end{itemize}

\sphinxAtStartPar
\sphinxstylestrong{Engine}
\begin{itemize}
\item {} 
\sphinxAtStartPar
Iniciar um projeto no Defold;

\item {} 
\sphinxAtStartPar
Navegar nas áreas: área de trabalho, Terminal.

\item {} 
\sphinxAtStartPar
Para os alunos irem se ambientando e familiarizando com a ferramentas é necessário fazer um exploratório principlamente do Defold.

\end{itemize}

\sphinxAtStartPar
\sphinxstylestrong{Fundamentos dos Jogos}
\begin{itemize}
\item {} 
\sphinxAtStartPar
Termos e definições de tipos de imagens (JPG,PNG,GIF E SVG)

\item {} 
\sphinxAtStartPar
Termos e definições: Sprites, Plataformas, Engines, Tiles

\item {} 
\sphinxAtStartPar
Exemplos de Jogos que foram desenvolvidos utilizando a linguagem.

\end{itemize}

\sphinxAtStartPar
\sphinxstylestrong{Ferramentas Gráficas}
\begin{itemize}
\item {} 
\sphinxAtStartPar
Libre Sprite

\item {} 
\sphinxAtStartPar
Inkscape

\end{itemize}
\subsubsection*{Descrição}

\sphinxAtStartPar
A aula abordará os conceitos introdutórios sobre jogos e seu funcionamento,
incluindo alguns exemplos. Diversos termos serão abordados na aula, introduzindo o aluno
ao mundo do desenvolvimento de jogos suas ferramentas principais.

\sphinxAtStartPar
Ao final, será reservado um momento para que as crianças possam explorar livremente essas ferramentas,
incentivando a criatividade, a experimentação e o desenvolvimento de habilidades artísticas e técnicas.
Após a apresentação dos Slides é necessário deixar pelo menos 30 minutos para as crianças brincarem.

\newpage


\section{Exercício Aula 1 \sphinxhyphen{} Introdução}
\label{\detokenize{modelo_aula_01:exercicio-aula-1-introducao}}
\sphinxAtStartPar
Data: \_\_\_/\_\_\_/\_\_\_

\sphinxAtStartPar
\sphinxstylestrong{MARQUE COM X A RESPOSTA CORRETA}

\sphinxAtStartPar
1 \sphinxhyphen{} O que são Tile Maps?
\begin{itemize}
\item {} 
\sphinxAtStartPar
{[}  {]} Imagens para serem utilizadas no jogo.

\item {} 
\sphinxAtStartPar
{[}  {]} Um programa.

\item {} 
\sphinxAtStartPar
{[}  {]} Um jogo que pode ser aberto no computador.

\item {} 
\sphinxAtStartPar
{[}  {]} Um personagem.

\end{itemize}

\sphinxAtStartPar
2 \textendash{} Qual o tipo/formato de imagem corresponde a uma animação?
\begin{itemize}
\item {} 
\sphinxAtStartPar
{[}  {]} JPG

\item {} 
\sphinxAtStartPar
{[}  {]} PNG

\item {} 
\sphinxAtStartPar
{[}  {]} WMA

\item {} 
\sphinxAtStartPar
{[}  {]} GIF

\end{itemize}

\sphinxAtStartPar
3 \textendash{} O que é uma Engine com suas palavras?

\sphinxAtStartPar
4 \sphinxhyphen{} O que são Sprites?
\begin{itemize}
\item {} 
\sphinxAtStartPar
{[}  {]} Imagens que são utilizadas para construção do jogo.

\item {} 
\sphinxAtStartPar
{[}  {]} Códigos que são digitados para construção do jogo.

\item {} 
\sphinxAtStartPar
{[}  {]} Software para construção do jogo.

\item {} 
\sphinxAtStartPar
{[}  {]} Imagens no formato JPG que se movimentam.

\end{itemize}

\sphinxAtStartPar
5 \textendash{} Qual a linguagem de programação será utilizada no curso?
\begin{itemize}
\item {} 
\sphinxAtStartPar
{[} {]} Java

\item {} 
\sphinxAtStartPar
{[} {]} Limon

\item {} 
\sphinxAtStartPar
{[} {]} Python

\item {} 
\sphinxAtStartPar
{[} {]} Lua

\end{itemize}

\newpage
\subsubsection*{Bibliografia}

\sphinxAtStartPar
LUA. Lua.org. Rio de Janeiro: PUC\sphinxhyphen{}Rio, {[}s.d.{]}.
Disponível em: \sphinxurl{https://www.lua.org/}. Acesso em: 14 ago. 2025.

\sphinxAtStartPar
LIBRESPRITE. LibreSprite. {[}S. l.{]}: LibreSprite, {[}s.d.{]}.
Disponível em: \sphinxurl{https://libresprite.github.io/}. Acesso em: 14 ago. 2025.

\sphinxAtStartPar
INKSCAPE. Inkscape. {[}S. l.{]}: Inkscape, {[}s.d.{]}.
Disponível em: \sphinxurl{https://inkscape.org/}. Acesso em: 14 ago. 2025.

\sphinxAtStartPar
MICROSOFT. Visual Studio Code. {[}S. l.{]}: Microsoft, {[}s.d.{]}.
Disponível em: \sphinxurl{https://code.visualstudio.com/}. Acesso em: 14 ago. 2025.

\sphinxAtStartPar
Rabin, S. (n.d.). Introdução ao desenvolvimento de games (Cengage Le, Vol. 1).
\sphinxurl{https://www.amazon.com.br}/Introdução\sphinxhyphen{}desenvolvimento\sphinxhyphen{}games\sphinxhyphen{}Entendendo\sphinxhyphen{}universo/dp/852211143X

\sphinxAtStartPar
Rita, S. (2017). Treinamento em Lógica de Programação (R. Sandra (ed.); Universo d).
\sphinxurl{https://br.99ebooks.net/download/treinamento-em-logica-de-programacao/}

\sphinxstepscope


\chapter{Introdução ao Defold}
\label{\detokenize{modelo_aula_02:introducao-ao-defold}}\label{\detokenize{modelo_aula_02::doc}}

\section{Aula sobre Engine Defold}
\label{\detokenize{modelo_aula_02:aula-sobre-engine-defold}}
\sphinxAtStartPar
\sphinxstylestrong{Tempo de Aula:} 1h30min
\subsubsection*{Recursos Didáticos}

\sphinxAtStartPar
Projetor, quadro branco, computadores.
\subsubsection*{Metodologia}
\begin{itemize}
\item {} 
\sphinxAtStartPar
Expositiva.

\item {} 
\sphinxAtStartPar
Exercícios Guiados em programação e operação da Engine.

\item {} 
\sphinxAtStartPar
Estimulo a participação.

\end{itemize}
\subsubsection*{Objetivos}
\begin{itemize}
\item {} 
\sphinxAtStartPar
Compreender conceitos básicos que estão por trás da criação de jogos.

\item {} 
\sphinxAtStartPar
Aplicar na prática alteração e criação de partes de um projeto.

\item {} 
\sphinxAtStartPar
Desenvolver curiosidade e autonomia e raciocínio lógico.

\end{itemize}
\subsubsection*{Conteúdo Programático}

\sphinxAtStartPar
\sphinxstylestrong{Ambientação}
\begin{itemize}
\item {} 
\sphinxAtStartPar
Estrutura de um projeto

\item {} 
\sphinxAtStartPar
Revisão conceitos(Sprites, Tilesource, TileMaps)

\end{itemize}

\sphinxAtStartPar
\sphinxstylestrong{Estudo de Caso}
\begin{itemize}
\item {} 
\sphinxAtStartPar
Analise Projeto simples

\end{itemize}

\sphinxAtStartPar
\sphinxstylestrong{Aplicação dos conceitos}
\begin{itemize}
\item {} 
\sphinxAtStartPar
Criando personagem e cenários e adicionando ao projeto.

\item {} 
\sphinxAtStartPar
Momento para criar, desenhar personagens.

\end{itemize}
\subsubsection*{Descrição}

\sphinxAtStartPar
No decorrer dessa aula será exposto conceitos e um projeto como estudo de caso, os alunos poderão interagir
e alterar e adicionar funcionalidades ao projeto. A ideia dessa aula é deixar as crianças a vontade para
alterar e brincar com a ferramenta. Deve\sphinxhyphen{}se apresentar a documentação aos alunos de uma maneira bem dinâmica,
exemplos de códigos na linguagem lua, podendo\sphinxhyphen{}se alternar os exemplos de código na Engine Defold e VSCODE.
\sphinxSetupCaptionForVerbatim{Exemplo exibição na tela com a função io.write.}
\def\sphinxLiteralBlockLabel{\label{\detokenize{modelo_aula_02:id1}}}
\begin{sphinxVerbatim}[commandchars=\\\{\}]
\PYG{+w}{    }\PYG{n+nb}{io.write}\PYG{p}{(}\PYG{l+s+s2}{\PYGZdq{}}\PYG{l+s+se}{\PYGZbs{}n}\PYG{l+s+s2}{Naruto }\PYG{l+s+se}{\PYGZbs{}n}\PYG{l+s+s2}{\PYGZdq{}}\PYG{p}{)}
\PYG{+w}{    }\PYG{n+nb}{io.write}\PYG{p}{(}\PYG{l+s+s2}{\PYGZdq{}}\PYG{l+s+s2}{Full Metal Alchemist }\PYG{l+s+se}{\PYGZbs{}n}\PYG{l+s+s2}{\PYGZdq{}}\PYG{p}{)}
\PYG{+w}{    }\PYG{n+nb}{io.write}\PYG{p}{(}\PYG{l+s+s2}{\PYGZdq{}}\PYG{l+s+s2}{Dragon Ball }\PYG{l+s+se}{\PYGZbs{}n}\PYG{l+s+s2}{\PYGZdq{}}\PYG{p}{)}
\end{sphinxVerbatim}

\sphinxAtStartPar
Da mesma forma é possível utilizar a função print para poder ver no terminal alguma informaçãO
\sphinxSetupCaptionForVerbatim{Exemplo exibição na tela com a função Print.}
\def\sphinxLiteralBlockLabel{\label{\detokenize{modelo_aula_02:id2}}}
\begin{sphinxVerbatim}[commandchars=\\\{\}]
\PYG{+w}{    }\PYG{n+nb}{print}\PYG{p}{(}\PYG{l+s+s2}{\PYGZdq{}}\PYG{l+s+se}{\PYGZbs{}n}\PYG{l+s+s2}{Naruto }\PYG{l+s+se}{\PYGZbs{}n}\PYG{l+s+s2}{\PYGZdq{}}\PYG{p}{)}
\PYG{+w}{    }\PYG{n+nb}{print}\PYG{p}{(}\PYG{l+s+s2}{\PYGZdq{}}\PYG{l+s+s2}{Full Metal Alchemist }\PYG{l+s+se}{\PYGZbs{}n}\PYG{l+s+s2}{\PYGZdq{}}\PYG{p}{)}
\PYG{+w}{    }\PYG{n+nb}{print}\PYG{p}{(}\PYG{l+s+s2}{\PYGZdq{}}\PYG{l+s+s2}{Dragon Ball }\PYG{l+s+se}{\PYGZbs{}n}\PYG{l+s+s2}{\PYGZdq{}}\PYG{p}{)}
\end{sphinxVerbatim}

\sphinxAtStartPar
Após os Alunos terem um entendimento básico da Engine podese aprofundar no entendimento de alguns conceitos
ja plicando na Engine:
\sphinxSetupCaptionForVerbatim{Exemplo de variáveis}
\def\sphinxLiteralBlockLabel{\label{\detokenize{modelo_aula_02:id3}}}
\begin{sphinxVerbatim}[commandchars=\\\{\}]
\PYG{+w}{    }\PYG{n+nv}{VIDA\PYGZus{}PLAYER}\PYG{+w}{ }\PYG{o}{=}\PYG{+w}{ }\PYG{l+m+mi}{100}
\PYG{+w}{    }\PYG{n+nv}{VIDA\PYGZus{}INIMIGO}\PYG{+w}{ }\PYG{o}{=}\PYG{+w}{ }\PYG{l+m+mi}{100}

\PYG{+w}{    }\PYG{n+nv}{DANO\PYGZus{}PLAYER}\PYG{+w}{ }\PYG{o}{=}\PYG{+w}{ }\PYG{l+m+mi}{20}
\PYG{+w}{    }\PYG{n+nv}{DANO\PYGZus{}INIMIGO}\PYG{+w}{ }\PYG{o}{=}\PYG{+w}{ }\PYG{l+m+mi}{15}

\PYG{+w}{    }\PYG{k+kd}{local}\PYG{+w}{ }\PYG{n+nv}{ataque\PYGZus{}player}\PYG{+w}{ }\PYG{o}{=}\PYG{+w}{ }\PYG{n+nv}{VIDA\PYGZus{}INIMIGO}\PYG{+w}{ }\PYG{o}{\PYGZhy{}}\PYG{+w}{ }\PYG{n+nv}{DANO\PYGZus{}PLAYER}
\PYG{+w}{    }\PYG{k+kd}{local}\PYG{+w}{ }\PYG{n+nv}{ataque\PYGZus{}inimigo}\PYG{+w}{ }\PYG{o}{=}\PYG{+w}{ }\PYG{n+nv}{VIDA\PYGZus{}PLAYER}\PYG{+w}{ }\PYG{o}{\PYGZhy{}}\PYG{+w}{ }\PYG{n+nv}{DANO\PYGZus{}INIMIGO}

\PYG{+w}{    }\PYG{n+nb}{print}\PYG{p}{(}\PYG{l+s+s2}{\PYGZdq{}}\PYG{l+s+s2}{Vida do Player: }\PYG{l+s+s2}{\PYGZdq{}}\PYG{+w}{ }\PYG{o}{..}\PYG{+w}{ }\PYG{n+nv}{VIDA\PYGZus{}PLAYER}\PYG{p}{)}
\PYG{+w}{    }\PYG{n+nb}{print}\PYG{p}{(}\PYG{l+s+s2}{\PYGZdq{}}\PYG{l+s+s2}{Vida do Inimigo: }\PYG{l+s+s2}{\PYGZdq{}}\PYG{+w}{ }\PYG{o}{..}\PYG{+w}{ }\PYG{n+nv}{VIDA\PYGZus{}INIMIGO}\PYG{p}{)}

\PYG{+w}{    }\PYG{n+nb}{print}\PYG{p}{(}\PYG{l+s+s2}{\PYGZdq{}}\PYG{l+s+s2}{após os ataques:}\PYG{l+s+s2}{\PYGZdq{}}\PYG{p}{)}

\PYG{+w}{    }\PYG{n+nb}{print}\PYG{p}{(}\PYG{l+s+s2}{\PYGZdq{}}\PYG{l+s+s2}{Vida do Player: }\PYG{l+s+s2}{\PYGZdq{}}\PYG{+w}{ }\PYG{o}{..}\PYG{+w}{ }\PYG{n+nv}{ataque\PYGZus{}player}\PYG{p}{)}
\PYG{+w}{    }\PYG{n+nb}{print}\PYG{p}{(}\PYG{l+s+s2}{\PYGZdq{}}\PYG{l+s+s2}{Vida do Inimigo: }\PYG{l+s+s2}{\PYGZdq{}}\PYG{+w}{ }\PYG{o}{..}\PYG{+w}{ }\PYG{n+nv}{ataque\PYGZus{}inimigo}\PYG{p}{)}
\end{sphinxVerbatim}

\newpage


\section{Exercício Aula 2 \sphinxhyphen{} Introdução ao Defold}
\label{\detokenize{modelo_aula_02:exercicio-aula-2-introducao-ao-defold}}
\sphinxAtStartPar
Data: \_\_\_/\_\_\_/\_\_\_

\sphinxAtStartPar
\sphinxstylestrong{Faça os exercícios:}
\begin{itemize}
\item {} 
\sphinxAtStartPar
1 \textendash{} Utilize os sprites de personagem, para importar para o Defold com as devidas configurações, anime o personagem, nos modos andando e parado.

\item {} 
\sphinxAtStartPar
2 \textendash{} Anime e configure o personagem com os comandos para direita e para a esquerda.

\item {} 
\sphinxAtStartPar
3 \textendash{} Crie um personagem animado com libresprite salve os arquivos e importe para o seu projeto.

\end{itemize}

\newpage
\subsubsection*{Bibliografia}

\sphinxAtStartPar
FOUNDATION, Defold. Defold \sphinxhyphen{} Official Homepage \sphinxhyphen{} Cross platform game engine.
Disponível em: \textless{}\sphinxurl{https://defold.com/}\textgreater{}. Acesso em: 25 jun. 2025.

\sphinxAtStartPar
Rabin, S. (n.d.). Introdução ao desenvolvimento de games (Cengage Le, Vol. 1).
\sphinxurl{https://www.amazon.com.br}/Introdução\sphinxhyphen{}desenvolvimento\sphinxhyphen{}games\sphinxhyphen{}Entendendo\sphinxhyphen{}universo/dp/852211143X

\sphinxAtStartPar
Rita, S. (2017). Treinamento em Lógica de Programação (R. Sandra (ed.); Universo d).
\sphinxurl{https://br.99ebooks.net/download/treinamento-em-logica-de-programacao/}

\sphinxstepscope


\chapter{Linguagem Lua I \sphinxhyphen{} variáveis e condições}
\label{\detokenize{modelo_aula_03:linguagem-lua-i-variaveis-e-condicoes}}\label{\detokenize{modelo_aula_03::doc}}
\sphinxAtStartPar
\sphinxstylestrong{Tempo de Aula:} 1h30min
\subsubsection*{Recursos Didáticos}

\sphinxAtStartPar
Projetor, quadro branco, computadores.
\subsubsection*{Metodologia}
\begin{itemize}
\item {} 
\sphinxAtStartPar
Expositiva.

\item {} 
\sphinxAtStartPar
Exercícios Guiados em programação e operação da Engine.

\item {} 
\sphinxAtStartPar
Estimulo a participação.

\item {} 
\sphinxAtStartPar
Desafio.

\end{itemize}
\subsubsection*{Objetivos}
\begin{itemize}
\item {} 
\sphinxAtStartPar
Compreender conceitos da linguagem de programação lua e aplicar na engine.

\item {} 
\sphinxAtStartPar
Operadores aritméticos

\item {} 
\sphinxAtStartPar
if e else

\item {} 
\sphinxAtStartPar
variáveis

\item {} 
\sphinxAtStartPar
Desenvolver curiosidade e autonomia e raciocínio lógico.

\end{itemize}
\subsubsection*{Conteúdo Programático}
\begin{itemize}
\item {} 
\sphinxAtStartPar
Programar teclas/Clicks para que o programa reconheça os comandos.

\item {} 
\sphinxAtStartPar
Animar personagem.

\item {} 
\sphinxAtStartPar
Entender conceito de variável e função.

\end{itemize}
\subsubsection*{Descrição}

\sphinxAtStartPar
Deve\sphinxhyphen{}se gastar pelo menos 20 a 30 minutos reforçando os conteúdos com alguns
exercícios na linguagem lua. Essa aula pode ser utilizada para focar exclusivamente nas funções e
conceitos da linguagem lua.
\sphinxSetupCaptionForVerbatim{Exemplo de Entrada}
\def\sphinxLiteralBlockLabel{\label{\detokenize{modelo_aula_03:id1}}}
\begin{sphinxVerbatim}[commandchars=\\\{\}]
\PYG{+w}{    }\PYG{n+nb}{print}\PYG{p}{(}\PYG{l+s+s2}{\PYGZdq{}}\PYG{l+s+s2}{Digite o jogo de celular:}\PYG{l+s+s2}{\PYGZdq{}}\PYG{p}{)}
\PYG{+w}{    }\PYG{n+nv}{nome}\PYG{+w}{ }\PYG{o}{=}\PYG{+w}{ }\PYG{n+nb}{io.read}\PYG{p}{()}
\PYG{+w}{    }\PYG{n+nb}{print}\PYG{p}{(}\PYG{l+s+s2}{\PYGZdq{}}\PYG{l+s+s2}{Digite a nota}\PYG{l+s+s2}{\PYGZdq{}}\PYG{p}{)}
\PYG{+w}{    }\PYG{n+nv}{nota}\PYG{+w}{ }\PYG{o}{=}\PYG{+w}{ }\PYG{n+nb}{io.read}\PYG{p}{()}
\PYG{+w}{    }\PYG{n+nb}{print}\PYG{p}{(}\PYG{l+s+s2}{\PYGZdq{}}\PYG{l+s+s2}{digite a descricao}\PYG{l+s+s2}{\PYGZdq{}}\PYG{p}{)}
\PYG{+w}{    }\PYG{n+nv}{descricao}\PYG{+w}{ }\PYG{o}{=}\PYG{+w}{ }\PYG{n+nb}{io.read}\PYG{p}{()}
\PYG{+w}{    }\PYG{n+nb}{print}\PYG{p}{(}\PYG{l+s+s2}{\PYGZdq{}}\PYG{l+s+s2}{O nome do jogo e :}\PYG{l+s+s2}{\PYGZdq{}}\PYG{p}{,}\PYG{n+nv}{nome}\PYG{p}{,}\PYG{+w}{ }\PYG{l+s+s2}{\PYGZdq{}}\PYG{l+s+s2}{ e a nota dele e }\PYG{l+s+s2}{\PYGZdq{}}\PYG{p}{,}\PYG{n+nv}{nota}\PYG{p}{)}
\PYG{+w}{    }\PYG{n+nb}{print}\PYG{p}{(}\PYG{l+s+s2}{\PYGZdq{}}\PYG{l+s+s2}{Descricao : }\PYG{l+s+s2}{\PYGZdq{}}\PYG{p}{,}\PYG{n+nv}{descricao}\PYG{p}{)}
\end{sphinxVerbatim}
\sphinxSetupCaptionForVerbatim{Exemplo de variáveis}
\def\sphinxLiteralBlockLabel{\label{\detokenize{modelo_aula_03:id2}}}
\begin{sphinxVerbatim}[commandchars=\\\{\}]
\PYG{+w}{    }\PYG{k+kd}{local}\PYG{+w}{ }\PYG{n+nv}{idade}\PYG{+w}{ }\PYG{o}{=}\PYG{+w}{ }\PYG{l+m+mi}{20}

\PYG{+w}{    }\PYG{k+kr}{if}\PYG{+w}{ }\PYG{n+nv}{idade}\PYG{+w}{ }\PYG{o}{\PYGZgt{}=}\PYG{+w}{ }\PYG{l+m+mi}{18}\PYG{+w}{ }\PYG{k+kr}{then}
\PYG{+w}{        }\PYG{n+nb}{print}\PYG{p}{(}\PYG{l+s+s2}{\PYGZdq{}}\PYG{l+s+s2}{Você é maior de idade.}\PYG{l+s+s2}{\PYGZdq{}}\PYG{p}{)}
\PYG{+w}{    }\PYG{k+kr}{else}
\PYG{+w}{        }\PYG{n+nb}{print}\PYG{p}{(}\PYG{l+s+s2}{\PYGZdq{}}\PYG{l+s+s2}{Você é menor de idade.}\PYG{l+s+s2}{\PYGZdq{}}\PYG{p}{)}
\PYG{+w}{    }\PYG{k+kr}{end}
\end{sphinxVerbatim}
\sphinxSetupCaptionForVerbatim{Exemplo operadores}
\def\sphinxLiteralBlockLabel{\label{\detokenize{modelo_aula_03:id3}}}
\begin{sphinxVerbatim}[commandchars=\\\{\}]
\PYG{+w}{    }\PYG{k+kd}{local}\PYG{+w}{ }\PYG{n+nv}{idade}\PYG{+w}{ }\PYG{o}{=}\PYG{+w}{ }\PYG{l+m+mi}{20}

\PYG{+w}{    }\PYG{k+kr}{if}\PYG{+w}{ }\PYG{n+nv}{idade}\PYG{+w}{ }\PYG{o}{\PYGZti{}=}\PYG{+w}{ }\PYG{l+m+mi}{18}\PYG{+w}{ }\PYG{k+kr}{then}
\PYG{+w}{        }\PYG{n+nb}{print}\PYG{p}{(}\PYG{l+s+s2}{\PYGZdq{}}\PYG{l+s+s2}{Você é maior de idade.}\PYG{l+s+s2}{\PYGZdq{}}\PYG{p}{)}
\PYG{+w}{    }\PYG{k+kr}{else}
\PYG{+w}{        }\PYG{n+nb}{print}\PYG{p}{(}\PYG{l+s+s2}{\PYGZdq{}}\PYG{l+s+s2}{Você é menor de idade.}\PYG{l+s+s2}{\PYGZdq{}}\PYG{p}{)}
\PYG{+w}{    }\PYG{k+kr}{end}
\end{sphinxVerbatim}
\sphinxSetupCaptionForVerbatim{Exemplo de operadores aritméticos}
\def\sphinxLiteralBlockLabel{\label{\detokenize{modelo_aula_03:id4}}}
\begin{sphinxVerbatim}[commandchars=\\\{\}]
\PYG{+w}{    }\PYG{k+kd}{local}\PYG{+w}{ }\PYG{n+nv}{idade}\PYG{+w}{ }\PYG{o}{=}\PYG{+w}{ }\PYG{l+m+mi}{20}
\PYG{+w}{    }\PYG{k+kd}{local}\PYG{+w}{ }\PYG{n+nv}{nome}\PYG{+w}{ }\PYG{o}{=}\PYG{+w}{ }\PYG{l+s+s2}{\PYGZdq{}}\PYG{l+s+s2}{Seu nome}\PYG{l+s+s2}{\PYGZdq{}}

\PYG{+w}{    }\PYG{k+kr}{if}\PYG{+w}{ }\PYG{n+nv}{idade}\PYG{+w}{ }\PYG{o}{\PYGZgt{}=}\PYG{+w}{ }\PYG{l+m+mi}{18}\PYG{+w}{ }\PYG{k+kr}{then}
\PYG{+w}{        }\PYG{n+nb}{print}\PYG{p}{(}\PYG{l+s+s2}{\PYGZdq{}}\PYG{l+s+s2}{Você é maior de idade.}\PYG{l+s+s2}{\PYGZdq{}}\PYG{p}{,}\PYG{n+nv}{nome}\PYG{p}{)}
\PYG{+w}{    }\PYG{k+kr}{else}
\PYG{+w}{        }\PYG{n+nb}{print}\PYG{p}{(}\PYG{l+s+s2}{\PYGZdq{}}\PYG{l+s+s2}{Você é menor de idade.}\PYG{l+s+s2}{\PYGZdq{}}\PYG{p}{,}\PYG{n+nv}{nome}\PYG{p}{)}
\PYG{+w}{    }\PYG{k+kr}{end}
\end{sphinxVerbatim}

\newpage
\subsubsection*{Bibliografia}

\sphinxAtStartPar
FOUNDATION, Defold. Defold \sphinxhyphen{} Official Homepage \sphinxhyphen{} Cross platform game engine.
Disponível em: \textless{}\sphinxurl{https://defold.com/}\textgreater{}. Acesso em: 25 jun. 2025.

\sphinxAtStartPar
Rabin, S. (n.d.). Introdução ao desenvolvimento de games (Cengage Le, Vol. 1).
\sphinxurl{https://www.amazon.com.br}/Introdução\sphinxhyphen{}desenvolvimento\sphinxhyphen{}games\sphinxhyphen{}Entendendo\sphinxhyphen{}universo/dp/852211143X

\sphinxAtStartPar
Rita, S. (2017). Treinamento em Lógica de Programação (R. Sandra (ed.); Universo d).
\sphinxurl{https://br.99ebooks.net/download/treinamento-em-logica-de-programacao/}

\sphinxstepscope


\chapter{Linguagem Lua II \sphinxhyphen{} laços de repetição}
\label{\detokenize{modelo_aula_04:linguagem-lua-ii-lacos-de-repeticao}}\label{\detokenize{modelo_aula_04::doc}}
\sphinxAtStartPar
\sphinxstylestrong{Tempo de Aula:} 1h30min
\subsubsection*{Recursos Didáticos}

\sphinxAtStartPar
Projetor, quadro branco, computadores.
\subsubsection*{Metodologia}
\begin{itemize}
\item {} 
\sphinxAtStartPar
Expositiva.

\item {} 
\sphinxAtStartPar
Exercícios Guiados em programação e operação da Engine.

\item {} 
\sphinxAtStartPar
Desafio \sphinxhyphen{} Memória Muscular

\end{itemize}
\subsubsection*{Conteúdo Programático}
\begin{itemize}
\item {} 
\sphinxAtStartPar
Utilização de funções da Linguagem Lua.

\item {} 
\sphinxAtStartPar
Utilização de variáveis, laços e condições.

\item {} 
\sphinxAtStartPar
Utilização de Tabelas.

\end{itemize}
\subsubsection*{Descrição}

\sphinxAtStartPar
Para Os Alunos internalizarem bem como é fluxo de funcionamento do Defold com relação aos
escripts, Serão entregues exercícios que ja utilizam de determinadas funções do próprio Defold.
\sphinxSetupCaptionForVerbatim{Exemplo de laço for}
\def\sphinxLiteralBlockLabel{\label{\detokenize{modelo_aula_04:id1}}}
\begin{sphinxVerbatim}[commandchars=\\\{\}]
\PYG{+w}{    }\PYG{k+kr}{for}\PYG{+w}{ }\PYG{n+nv}{i}\PYG{o}{=}\PYG{l+m+mi}{0}\PYG{p}{,}\PYG{+w}{ }\PYG{l+m+mi}{10}\PYG{+w}{ }\PYG{k+kr}{do}
\PYG{+w}{    }\PYG{n+nb}{print}\PYG{p}{(}\PYG{l+s+s2}{\PYGZdq{}}\PYG{l+s+s2}{status}\PYG{l+s+s2}{\PYGZdq{}}\PYG{p}{,}\PYG{n+nv}{i}\PYG{p}{)}
\PYG{+w}{    }\PYG{k+kr}{end}
\end{sphinxVerbatim}
\sphinxSetupCaptionForVerbatim{Exemplo de laço while}
\def\sphinxLiteralBlockLabel{\label{\detokenize{modelo_aula_04:id2}}}
\begin{sphinxVerbatim}[commandchars=\\\{\}]
\PYG{+w}{    }\PYG{k+kd}{local}\PYG{+w}{ }\PYG{n+nv}{i}\PYG{+w}{ }\PYG{o}{=}\PYG{+w}{ }\PYG{l+m+mi}{1}

\PYG{+w}{    }\PYG{k+kr}{while}\PYG{+w}{ }\PYG{n+nv}{i}\PYG{+w}{ }\PYG{o}{\PYGZlt{}=}\PYG{+w}{ }\PYG{l+m+mi}{5}\PYG{+w}{ }\PYG{k+kr}{do}
\PYG{+w}{        }\PYG{n+nb}{print}\PYG{p}{(}\PYG{n+nv}{i}\PYG{p}{)}
\PYG{+w}{        }\PYG{n+nv}{i}\PYG{+w}{ }\PYG{o}{=}\PYG{+w}{ }\PYG{n+nv}{i}\PYG{+w}{ }\PYG{o}{+}\PYG{+w}{ }\PYG{l+m+mi}{1}
\PYG{+w}{    }\PYG{k+kr}{end}
\end{sphinxVerbatim}
\sphinxSetupCaptionForVerbatim{Exemplo de função}
\def\sphinxLiteralBlockLabel{\label{\detokenize{modelo_aula_04:id3}}}
\begin{sphinxVerbatim}[commandchars=\\\{\}]
\PYG{+w}{    }\PYG{k+kr}{function}\PYG{+w}{ }\PYG{n+nf}{idade}\PYG{p}{(}\PYG{n+nv}{d}\PYG{p}{)}
\PYG{+w}{        }\PYG{n+nb}{print}\PYG{p}{(}\PYG{n+nv}{d}\PYG{p}{)}
\PYG{+w}{    }\PYG{k+kr}{end}
\PYG{+w}{    }\PYG{n+nf}{idade}\PYG{p}{(}\PYG{l+m+mi}{20}\PYG{p}{)}
\end{sphinxVerbatim}

\newpage
\subsubsection*{Bibliografia}

\sphinxAtStartPar
FOUNDATION, Defold. Defold \sphinxhyphen{} Official Homepage \sphinxhyphen{} Cross platform game engine.
Disponível em: \textless{}\sphinxurl{https://defold.com/}\textgreater{}. Acesso em: 25 jun. 2025.

\sphinxAtStartPar
Rabin, S. (n.d.). Introdução ao desenvolvimento de games (Cengage Le, Vol. 1).
\sphinxurl{https://www.amazon.com.br}/Introdução\sphinxhyphen{}desenvolvimento\sphinxhyphen{}games\sphinxhyphen{}Entendendo\sphinxhyphen{}universo/dp/852211143X

\sphinxAtStartPar
Rita, S. (2017). Treinamento em Lógica de Programação (R. Sandra (ed.); Universo d).
\sphinxurl{https://br.99ebooks.net/download/treinamento-em-logica-de-programacao/}

\sphinxstepscope


\chapter{Lua III \sphinxhyphen{} Tabelas}
\label{\detokenize{modelo_aula_05:lua-iii-tabelas}}\label{\detokenize{modelo_aula_05::doc}}
\sphinxAtStartPar
\sphinxstylestrong{Tempo de Aula:} 1h30min
\subsubsection*{Recursos Didáticos}

\sphinxAtStartPar
Projetor, quadro branco, computadores.
\subsubsection*{Metodologia}
\begin{itemize}
\item {} 
\sphinxAtStartPar
Expositiva.

\item {} 
\sphinxAtStartPar
Exercícios Guiados em programação e operação da Engine.

\item {} 
\sphinxAtStartPar
Desafio \sphinxhyphen{} Memória Muscular

\end{itemize}
\subsubsection*{Conteúdo Programático}
\begin{itemize}
\item {} 
\sphinxAtStartPar
Navegação na plataforma.

\item {} 
\sphinxAtStartPar
Aplicação do conceito de tabelas no Defold.

\item {} 
\sphinxAtStartPar
Exercícios Dinãmicos.

\end{itemize}
\subsubsection*{Descrição}

\sphinxAtStartPar
Utilizar os conceitos das aulas anteriores para aplicar no Defold
com exemplos mais complexos.
\subsubsection*{Avaliação}

\sphinxAtStartPar
Execução dos exercícios em sala de aula.
\sphinxSetupCaptionForVerbatim{Exemplo de tabelas}
\def\sphinxLiteralBlockLabel{\label{\detokenize{modelo_aula_05:id1}}}
\begin{sphinxVerbatim}[commandchars=\\\{\}]
\PYG{+w}{    }\PYG{c+c1}{\PYGZhy{}\PYGZhy{} Criando uma caixa vazia}
\PYG{+w}{    }\PYG{n+nv}{meu\PYGZus{}bau}\PYG{+w}{ }\PYG{o}{=}\PYG{+w}{ }\PYG{p}{\PYGZob{}\PYGZcb{}}

\PYG{+w}{    }\PYG{c+c1}{\PYGZhy{}\PYGZhy{} Colocando coisas dentro}
\PYG{+w}{    }\PYG{n+nv}{meu\PYGZus{}bau}\PYG{p}{[}\PYG{l+m+mi}{1}\PYG{p}{]}\PYG{+w}{ }\PYG{o}{=}\PYG{+w}{ }\PYG{l+s+s2}{\PYGZdq{}}\PYG{l+s+s2}{Bola}\PYG{l+s+s2}{\PYGZdq{}}
\PYG{+w}{    }\PYG{n+nv}{meu\PYGZus{}bau}\PYG{p}{[}\PYG{l+m+mi}{2}\PYG{p}{]}\PYG{+w}{ }\PYG{o}{=}\PYG{+w}{ }\PYG{l+s+s2}{\PYGZdq{}}\PYG{l+s+s2}{Carrinho}\PYG{l+s+s2}{\PYGZdq{}}
\PYG{+w}{    }\PYG{n+nv}{meu\PYGZus{}bau}\PYG{p}{[}\PYG{l+m+mi}{3}\PYG{p}{]}\PYG{+w}{ }\PYG{o}{=}\PYG{+w}{ }\PYG{l+s+s2}{\PYGZdq{}}\PYG{l+s+s2}{Boneca}\PYG{l+s+s2}{\PYGZdq{}}
\PYG{+w}{    }\PYG{n+nb}{print}\PYG{p}{(}\PYG{n+nv}{meu\PYGZus{}bau}\PYG{p}{[}\PYG{l+m+mi}{1}\PYG{p}{])}\PYG{+w}{  }\PYG{c+c1}{\PYGZhy{}\PYGZhy{} Vai mostrar: Bola}
\PYG{+w}{    }\PYG{n+nb}{print}\PYG{p}{(}\PYG{n+nv}{meu\PYGZus{}bau}\PYG{p}{[}\PYG{l+m+mi}{3}\PYG{p}{])}\PYG{+w}{  }\PYG{c+c1}{\PYGZhy{}\PYGZhy{} Vai mostrar: Boneca}
\end{sphinxVerbatim}

\newpage
\subsubsection*{Bibliografia}

\sphinxAtStartPar
FOUNDATION, Defold. Defold \sphinxhyphen{} Official Homepage \sphinxhyphen{} Cross platform game engine.
Disponível em: \textless{}\sphinxurl{https://defold.com/}\textgreater{}. Acesso em: 25 jun. 2025.

\sphinxAtStartPar
Rabin, S. (n.d.). Introdução ao desenvolvimento de games (Cengage Le, Vol. 1).
\sphinxurl{https://www.amazon.com.br}/Introdução\sphinxhyphen{}desenvolvimento\sphinxhyphen{}games\sphinxhyphen{}Entendendo\sphinxhyphen{}universo/dp/852211143X

\sphinxAtStartPar
Rita, S. (2017). Treinamento em Lógica de Programação (R. Sandra (ed.); Universo d).
\sphinxurl{https://br.99ebooks.net/download/treinamento-em-logica-de-programacao/}

\sphinxstepscope


\chapter{Defold: Animações}
\label{\detokenize{modelo_aula_06:defold-animacoes}}\label{\detokenize{modelo_aula_06::doc}}
\sphinxAtStartPar
\sphinxstylestrong{Tempo de Aula:} 1h30min
\subsubsection*{Recursos Didáticos}

\sphinxAtStartPar
Projetor, quadro branco, computadores.
\subsubsection*{Metodologia}
\begin{itemize}
\item {} 
\sphinxAtStartPar
Expositiva.

\item {} 
\sphinxAtStartPar
Exercícios Guiados em programação e operação da Engine.

\item {} 
\sphinxAtStartPar
Desafio \sphinxhyphen{} Memória Muscular

\end{itemize}
\subsubsection*{Conteúdo Programático}
\begin{itemize}
\item {} 
\sphinxAtStartPar
Compreender as funções de animação

\item {} 
\sphinxAtStartPar
Aplicar laços, condições as animações

\item {} 
\sphinxAtStartPar
Criação de zonas de colizão

\end{itemize}
\subsubsection*{Descrição}

\sphinxAtStartPar
Iniciar um projeto do zero, buscando através da repetição fazer o aluno decorar os
caminhos, criar personagem simples no início, e animá\sphinxhyphen{}lo a partir dos frames.
É importante entender as zonas de colisão, pois com base nelas o personagem irá interagir
com o ambiente, o chão ou outros elementos.
Nessa aula é importante frisar as animações, e configurações.
\subsubsection*{Avaliação}

\sphinxAtStartPar
Execução dos exercícios em sala de aula.
\subsubsection*{Exercícios de Desafio}

\sphinxAtStartPar
\sphinxstylestrong{Desafio 1: Criar uma animação personalizada}
Utilize a ferramenta de edição de sprites para criar uma animação com pelo menos 5 frames diferentes.
Configure a velocidade da animação para que o movimento fique natural e fluido.

\sphinxAtStartPar
\sphinxstylestrong{Desafio 2: Ajustar zonas de colisão}
No seu personagem animado, habilite e ajuste as zonas de colisão para que elas estejam alinhadas corretamente com a forma do sprite em cada frame.
Teste a interação do personagem com o chão ou outros objetos para garantir que as colisões funcionem adequadamente.

\sphinxAtStartPar
\sphinxstylestrong{Desafio 3: Manipular pontos âncoras da imagem}
Experimente alterar os pontos âncoras do sprite para entender como eles afetam a posição e o comportamento da animação.
Crie um pequeno projeto onde o personagem gira ou se move com base nos pontos âncoras definidos.

\newpage
\subsubsection*{Bibliografia}

\sphinxAtStartPar
FOUNDATION, Defold. Defold \sphinxhyphen{} Official Homepage \sphinxhyphen{} Cross platform game engine.
Disponível em: \textless{}\sphinxurl{https://defold.com/}\textgreater{}. Acesso em: 25 jun. 2025.

\sphinxAtStartPar
Rabin, S. (n.d.). Introdução ao desenvolvimento de games (Cengage Le, Vol. 1).
\sphinxurl{https://www.amazon.com.br}/Introdução\sphinxhyphen{}desenvolvimento\sphinxhyphen{}games\sphinxhyphen{}Entendendo\sphinxhyphen{}universo/dp/852211143X

\sphinxAtStartPar
Rita, S. (2017). Treinamento em Lógica de Programação (R. Sandra (ed.); Universo d).
\sphinxurl{https://br.99ebooks.net/download/treinamento-em-logica-de-programacao/}

\sphinxstepscope


\chapter{Defold: GUI \sphinxhyphen{} e Assets}
\label{\detokenize{modelo_aula_07:defold-gui-e-assets}}\label{\detokenize{modelo_aula_07::doc}}
\sphinxAtStartPar
\sphinxstylestrong{Tempo de Aula:} 1h30min
\subsubsection*{Recursos Didáticos}

\sphinxAtStartPar
Projetor, quadro branco, computadores.
\subsubsection*{Metodologia}
\begin{itemize}
\item {} 
\sphinxAtStartPar
Expositiva.

\item {} 
\sphinxAtStartPar
Exercícios Guiados em programação e operação da Engine.

\item {} 
\sphinxAtStartPar
Desafio \sphinxhyphen{} Memória Muscular

\end{itemize}
\subsubsection*{Conteúdo Programático}
\begin{itemize}
\item {} 
\sphinxAtStartPar
Construção de um menu de exibição

\item {} 
\sphinxAtStartPar
Configurações de menssageria.

\end{itemize}
\subsubsection*{Descrição}

\sphinxAtStartPar
Nessa aula será demonstrado como criar um menu de acesso ao jogo e configurações de sons.
\subsubsection*{Avaliação}

\sphinxAtStartPar
Execução dos exercícios em sala de aula.

\newpage
\subsubsection*{Bibliografia}

\sphinxAtStartPar
SCIRRA, L. (2022). Manual Construct.
\sphinxurl{https://www.construct.net/en/make-games/manuals/construct-3}

\sphinxAtStartPar
Rabin, S. (n.d.). Introdução ao desenvolvimento de games (Cengage Le, Vol. 1).
\sphinxurl{https://www.amazon.com.br}/Introdução\sphinxhyphen{}desenvolvimento\sphinxhyphen{}games\sphinxhyphen{}Entendendo\sphinxhyphen{}universo/dp/852211143X

\sphinxAtStartPar
Rita, S. (2017). Treinamento em Lógica de Programação (R. Sandra (ed.); Universo d).
\sphinxurl{https://br.99ebooks.net/download/treinamento-em-logica-de-programacao/}

\sphinxAtStartPar
PENZILLA. Free Basic GUI Bundle | {[}Royalty Free{]}. itch.io, devlog em 12 mar. 2023.
Disponível em: \sphinxurl{https://penzilla.itch.io/basic-gui-bundle?download}. Acesso em: 15 ago. 2025.

\sphinxAtStartPar
KENNEY. Assets. Disponível em: \sphinxurl{https://kenney.nl/assets}. Acesso em: 15 ago. 2025.



\renewcommand{\indexname}{Índice}
\printindex
\end{document}